\documentclass[a4paper,12pt]{article}
\usepackage[utf8]{inputenc}
\usepackage[T1]{fontenc}
\usepackage[ngerman]{babel}
\usepackage{amsmath}
\usepackage{amssymb}
\usepackage{enumitem}
\usepackage{geometry}
\geometry{a4paper, margin=2.5cm}
\usepackage{parskip}

\title{Einsendeaufgaben -- Aufgabe 4.1\\[0.5em]
\large 61211 Analysis}
\author{}
\date{}

\begin{document}
\maketitle

\section*{Aufgabe 1}

\begin{enumerate}[label=\alph*)]

    \item
    Wir definieren:
    \[
        F: \mathbb{R}^2 \to \mathbb{R}, \quad \binom{x}{y} \mapsto x^2 + y + \sin(xy)
    \]

    \textit{Bestimmung der partiellen Ableitungen von $F$:}
    \begin{align*}
        D_1 F(x,y) &= 2x + y\cos(xy) \\
        D_2 F(x,y) &= 1 + x\cos(xy)
    \end{align*}

    Die partiellen Ableitungen $D_1 F$ und $D_2 F$ sind stetig und demnach ist $F$ partiell differenzierbar. Es gilt:
    \[
        F_y'(0,0) = 1
    \]
    und daraus folgt
    \[
        \operatorname{Rang} F_y'(0,0) = 1
    \]

    Wir können also Satz 4.2.3 über implizite Funktionen anwenden. Demzufolge existiert eine offene Umgebung $U$ von $0$ und eine offene Umgebung $V$ von $0$, sodass eine stetige Funktion $y: U \to V$ existiert mit:
    \[
        F(x, y(x)) = x^2 + y(x) + \sin(x\,y(x)) = 0
    \]
    Wir setzen $F(x, y(x))$ mit $x = 0$ ein und erhalten
    \[
        y(0) = 0
    \]
    Die stetige Funktion $y$ erfüllt also die Bedingung $y(0) = 0$.

    \item
    Nach Satz 4.2.3 ist $y$ stetig differenzierbar, also total differenzierbar, da die partiellen Ableitungen stetig sind. Der Definitionsbereich von $y$ ist die offene Umgebung $U$ von $x = 0$. $y$ ist also differenzierbar auf ganz $U$.

    \item
    Nach Satz 4.2.3(ii) gilt:
    \begin{align*}
        y'(0) &= -\left[ F_y'(0,0) \right]^{-1} F_x'(0,0) \\
              &= -\frac{1}{D_2 F(0,0)} D_1 F(0,0) \\
              &= -\frac{1}{1} \cdot 0 = 0
    \end{align*}

    Daraus folgt
    \[
        \lim_{x \to 0} \frac{y(x)}{x} = \lim_{x \to 0} \frac{y(x) - y(0)}{x - 0} = y'(0) = 0
    \]
    Nach dem $\varepsilon$-$\delta$-Kriterium existiert also ein $\delta > 0$, sodass $\frac{|y(x)|}{|x|} < 1$ und demnach $|y(x)| < |x|$ für alle $|x| < \delta$ gilt.

    Wir betrachten $y'(x)$ in der Umgebung
    \[
        U' := U \cap U_\delta(0) \cap U_1(0)
    \]

    Nach Satz 4.2.3(ii) gilt:
    \begin{align*}
        y'(x) &= -\left[ F_y'(x, y(x)) \right]^{-1} F_x'(0,0) \\
              &= -\frac{1}{D_2 F(x, y(x))} D_1 F(x, y(x)) \\
              &= -\frac{2x + y(x) \cos(x\,y(x))}{1 + x\cos(x\,y(x))}
    \end{align*}

    Wegen $x \in U_1(0)$, also $|x| < 1$ folgt
    \[
        1 + x\cos(x\,y(x)) > 0
    \]
    Der Nenner von $y'(x)$ ist also positiv für alle $x \in U'$!

    Außerdem gilt wegen $|y(x)| < |x|$
    \[
        y(x) \cos(x\,y(x)) \leq |y(x)| < |x| < |2x|.
    \]

    Mit den Abschätzungen können wir schließen
    \begin{align*}
        y'(x) &< 0 \quad \forall\, x \in U' \cap (0,1) \\
        y'(x) &> 0 \quad \forall\, x \in U' \cap (-1,0)
    \end{align*}

    Die Funktion $y|_{U' \cap (0,1)}$ ist demnach streng monoton fallend und $y|_{U' \cap (-1,0)}$ streng monoton steigend.

    \item
    Wir zeigen, dass es keine eindeutig definierte Funktion $x = x(y)$ in einer Umgebung $W$ des Punktes $(0,0)$ gibt.

    Sei $W \subseteq \mathbb{R}^2$ eine beliebige Umgebung des Punktes $(0,0)$. Wir können eine $\varepsilon$-Umgebung $W_\varepsilon(0,0) \subseteq W$ mit Norm $\|\cdot\|_\infty$ auswählen, mit der Bedingung:
    \begin{itemize}
        \item $U_\varepsilon(0) \subseteq \mathbb{R}$, sodass wir das Monotonieverhalten von $y(x)$ aus c) kennen
    \end{itemize}

    Da $y(x)$ stetig ist, existiert ein $\varepsilon' > 0$, mit $\varepsilon > \varepsilon' > 0$ und $y(U_{\varepsilon'}(0)) \subseteq U_\varepsilon(0)$.

    Wir betrachten die Umgebung
    \[
        U := U_{\varepsilon'}(0) \times U_\varepsilon(0) \subseteq W.
    \]
    Wegen dem Monotonieverhalten von $y(x)$ in $U_{\varepsilon'}(0)$ und $y(0) = 0$ gilt
    \[
        y(-\varepsilon') < 0 \quad \text{und} \quad y(\varepsilon') < 0.
    \]
    (wegen $y(U_{\varepsilon'}(0)) \subseteq U_\varepsilon(0)$ gilt $y \in U_\varepsilon(0)$.)

    Wir wählen ein $y \in \mathbb{R}$ aus, mit
    \[
        \max(y(-\varepsilon'),\, y(\varepsilon')) < y < 0.
    \]

    Wegen dem strengen Monotonieverhalten existieren $x_1, x_2 \in \mathbb{R}$ mit
    \[
        -\varepsilon < x_1 < 0 \quad \text{und} \quad 0 < x_2 < \varepsilon,
    \]
    sowie $y(x_1) = y = y(x_2)$.

    Wir haben also zwei Punkte $\binom{x_1}{y}$ und $\binom{x_2}{y}$ mit
    \[
        \binom{x_1}{y},\, \binom{x_2}{y} \in U \subseteq W
    \]
    und $F(x_1, y) = 0 = F(x_2, y)$.

    Demnach existiert kein eindeutiges $x(y)$ in der Umgebung $W$. Da $W$ beliebig war, existiert keine Umgebung des Punktes $(0,0)$ mit einer eindeutigen Funktion $x(y) = x$.

\end{enumerate}

\end{document}
